% Generated draft: requires mathematical and source review.
\documentclass[11pt]{article}
\usepackage[margin=25mm]{geometry}
\usepackage{fontspec}
\setmainfont{Latin Modern Roman}
\usepackage{amsmath,amssymb,mathtools,unicode-math}
\setmathfont{Latin Modern Math}
\usepackage{enumitem,needspace,xurl,hyperref}
\setlength{\parindent}{0pt}
\setlength{\parskip}{5pt}
\setlength{\emergencystretch}{3em}
\newcommand{\N}{\mathbb{N}}
\newcommand{\Z}{\mathbb{Z}}
\newcommand{\Q}{\mathbb{Q}}
\newcommand{\R}{\mathbb{R}}
\newcommand{\C}{\mathbb{C}}
\newcommand{\F}{\mathbb{F}}
\newcommand{\A}{\mathbb{A}}
\newcommand{\PP}{\mathbb{P}}
\newcommand{\E}{\mathbb{E}}
\newcommand{\eps}{\varepsilon}
\newcommand{\dd}{\,\mathrm{d}}
\begin{document}
\section*{UnsolvedMath ID: 701190}
\textbf{Source problem code:} catalogue entry
\subsection{Definitions and mathematical statement}
Let \(F_n\) denote the free group of rank \(n\). The supplied source does not specify the value or range of \(n\), or whether its question concerns a fixed rank or existence in some rank; the question below retains this rank parameter.

The commutator subgroup is
\[
[F_n,F_n]=\langle x^{-1}y^{-1}xy : x,y\in F_n\rangle.
\]
For \(r\in F_n\), its normal closure in \(F_n\) is
\[
\langle\!\langle r\rangle\!\rangle_{F_n}
=\langle grg^{-1}:g\in F_n\rangle.
\]
An element \(r\in F_n\) is a \emph{normal root} of \(u\in F_n\) if \(u\in\langle\!\langle r\rangle\!\rangle_{F_n}\). Two elements \(r,s\in F_n\) are conjugate if \(s=grg^{-1}\) for some \(g\in F_n\).

\textbf{Question (attributed in the source to W. Magnus).} Can there exist an element
\[
u\in F_n\setminus[F_n,F_n]
\]
and an infinite subset \(R\subseteq F_n\) such that
\[
\begin{gathered}
\forall r\in R,\qquad u\in\langle\!\langle r\rangle\!\rangle_{F_n},\\
\forall r,s\in R,\quad r\ne s\ \Longrightarrow\
\bigl(\forall g\in F_n,\ s\ne grg^{-1}\bigr)?
\end{gathered}
\]
Thus all normal closures and conjugacy relations are taken inside the same group \(F_n\).
\subsection{Short English statement}
Can an element of a free group lying outside its commutator subgroup have infinitely many pairwise nonconjugate normal roots? A normal root is an element whose normal closure contains the given element. The source leaves the rank parameter unspecified.
\subsection{Sources}
\begin{itemize}
\item Local UnsolvedMath registry: permanent ID \(701190\), exact source problem code \texttt{catalogue entry}, title \emph{Kourovka Notebook Problem 15.102}; supplied record updated \texttt{2026-10-16T14:00:00Z}.
\item Supplied bibliographic attribution: \emph{The Kourovka Notebook (Unsolved Problems in Group Theory)}, 15th issue (2002), Problem 15.102; also identified in the record with the 21st issue edition (2026), arXiv:\texttt{1401.0300v48}. The collection description names E. I. Khukhro and V. D. Mazurov as editors. These source identifications have not been independently verified.
\item The record lists V. E. Shpilrain as proposer, with proposed year 2002; its statement separately attributes the question to W. Magnus. Both attributions are retained.
\item The registry labels the problem \texttt{open}. This is a supplied status label, not an independently verified assessment of the literature.
\end{itemize}
\end{document}
